\section{Constructive upper certificates after exact tail normalization (NS89)}

\paragraph{Scope.}
The target remains an upper estimate $E_N\le C/\log N$ (or any fixed
positive logarithmic decay rate). No such cofinal estimate is proved here.
We construct explicit rational trial coefficients and certify their full
errors. The normalization argument below shows that this restricted search
preserves every possible upper decay rate, up to a uniform constant. It is
a Hilbert-space projection calculation in the classical Nyman--Beurling
setting, not a new arithmetic decay theorem. The earlier proposed
fixed-seed obstruction investigation was stopped when Alex redirected this
turn to upper bounds; no new lower-bound result is included.

Use $\mathcal H=L^2(0,\infty;dx)$, $\tau(x)=-\log x$ on $(0,1)$ and zero
elsewhere, $\sigma_n(x)=A(1/(nx))$, and
$A(t)=\int_0^t\{u\}\,du/u$. Let $G=G_N$ be the complete Gram,
$b_n=\langle\tau,\sigma_n\rangle$, $c=G^{-1}b$, and
$E_N=2-b^TG^{-1}b$. Every atom has exterior $1/(nx)$ for $x>1$.
Write $e(x)=x^{-1}\mathbf1_{x>1}$, $v_n=1/n$, and
\[
 \eta_N=v^Tc,\qquad q_N=v^TG^{-1}v=\|\Pi_N e\|^2>0.
\]

\paragraph{Exact optimal normalization and upper-rate transfer.}
The minimum error under the constraint $v^Ta=0$ is
\begin{equation}
 a_N^0=c-\frac{\eta_N}{q_N}G^{-1}v,\qquad
 E_N^0=E_N+\frac{\eta_N^2}{q_N}.
 \label{eq:ns89-optimal-normalization}
\end{equation}
Moreover, for every $N\ge1$,
\begin{equation}
 E_N\le E_N^0\le\frac{E_N}{q_N}\le G(1,1)E_N,
 \qquad G(1,1)<3.270466.
 \label{eq:ns89-upper-transfer}
\end{equation}
To prove the first formula, write $a=c+h$. Normal equations give the
full error $E_N+h^TGh$. Cauchy--Schwarz in the Gram metric minimizes this
under $v^Th=-\eta_N$, with the displayed optimizer. For the sharper upper
bound, $R_N=\tau-\Pi_N\tau$ is orthogonal to $V_N$, and
\[
 -\eta_N=\langle R_N,e\rangle
 =\langle R_N,(I-\Pi_N)e\rangle,
 \qquad \eta_N^2\le E_N(1-q_N).
\]
Finally $q_N\ge q_1=1/G(1,1)$. These are complete-space statements:
the exterior is removed by changing coefficients, never by deleting its
positive contribution from an unchanged norm. Thus a bound
$E_N^0\le C/\log N$ implies the desired original upper bound, and any
original bound of this order implies one for the normalized optimum with
constant at most $G(1,1)C$. The exponent is preserved. No assertion that
$q_N\to1$, that $\eta_N$ has a sign, or that either error tends to zero is
needed or made.

The finite certificate gives $1/q_{256}<1.144579$. Since the subspaces
are nested, $q_N$ is nondecreasing. Consequently the sharper comparison
\[
 E_N\le E_N^0\le1.144579\,E_N\qquad(N\ge256)
\]
holds at every later size. This is a cofinal comparison between two unknown
errors, not an upper decay estimate for either error.

\paragraph{Explicit rational witnesses.}
For $n\ge2$ round the proposed entries of $a_N^0$ to rational dyadics with
denominator $2^{64}$, then set
$\widetilde a_1=-\sum_{n=2}^N\widetilde a_n/n$ exactly. The first coefficient
is rational, not generally dyadic. The resulting vector has exactly zero
exterior. Its proof certificate is a fresh evaluation of
\[
 U_N=2-2b^T\widetilde a+\widetilde a^TG\widetilde a,
 \qquad E_N\le E_N^0\le U_N.
\]
The ball solve is never replaced by a midpoint inverse in a proof.
Rounding only chooses a trial vector, whose original full energy is then
recomputed. Since $v^T(\widetilde a-a_N^0)=0$, the constrained normal
equations give the additional check
\[
 U_N-E_N^0=(\widetilde a-a_N^0)^TG(\widetilde a-a_N^0).
\]

\paragraph{A finite arithmetic upper certificate including the remote tail.}
For any such zero-exterior coefficient vector $a$, put
\[
 \delta_m=\mathbf1_{m=1}+\sum_{n\mid m}a_n,\quad
 A_m=\sum_{k\le m}\delta_k,\quad
 B_m=\sum_{k\le m}\delta_k\log k,
\]
where $a_n=0$ for $n>N$. In reciprocal coordinates the exact residual on
$m\le t<m+1$ is $r(t)=A_m\log t-B_m$; it vanishes for $t<1$.
For $T>N$ an integer, define
\[
 H_T=\sum_{m=1}^{T-1}\int_m^{m+1}
          \frac{(A_m\log t-B_m)^2}{t^2}\,dt.
\]
All signed cross terms remain in this nonnegative cell expression.
The full Stirling estimate inherited from NS67/78 is
\[
 r(t)=u\log t+w+\epsilon(t),\quad |\epsilon(t)|\le D/t\quad(t\ge N),
\]
where $u=1-\tfrac12\sum a_n$,
$w=\tfrac12\sum a_n\log(n/(2\pi))$, and
$D=\tfrac16\sum n|a_n|$. With $L=\log T$, set
\[
 M_T=\frac{(uL+w)^2+2u(uL+w)+2u^2}{T},\qquad
 V_T=\frac{D^2}{3T^3}.
\]
Cauchy--Schwarz for the complete main/remainder cross term yields the
constructive upper certificate
\begin{equation}
 \|\tau-\sum a_n\sigma_n\|^2
 \le \mathcal U_T(a):=H_T+M_T+2\sqrt{M_TV_T}+V_T.
 \label{eq:ns89-physical-upper}
\end{equation}
This is a specialization and implementation of the already proved physical
tail estimate, not a new asymptotic tail theorem. For fixed $a$, the
certificate tends to its full error as $T\to\infty$.

\paragraph{What remains.}
The full upper-rate candidate would follow from an explicit family
$a^{(N)}$, with $\sum a_n^{(N)}/n=0$, and an independent uniform estimate
\[
 \mathcal U_{T_N}(a^{(N)})\le C/\log N
 \quad\hbox{for every sufficiently large }N.
\]
The finite certificates do not supply that estimate. In particular a small
cost for normalizing an already optimized vector does not prove the
optimized error decreases. The original signed divisor-cell sum $H_{T_N}$
still requires arithmetic control. These finite upper bounds are slightly
above the previously known unrestricted optima and are not advertised as
improved best finite bounds. Manuscript v1.66, G2 and RH remain unchanged.
